\section{Discussion and Conclusion}\label{sec:discussion}

\paragraph{Fidelity as a deployable signal.}
The paper's central empirical wager is that specification fidelity---cheap,
deterministic, and measurable on every run---tracks geometric correctness
closely enough to drive closed-loop intervention. The results above test the
wager in both directions: defect-reduction interventions improve
ground-truth metrics, and verifier verdicts predict geometric failures
without seeing any reference mesh.

\paragraph{Complementary interfaces.}
Treellis demonstrates the value of exposing editable structure to humans;
our specification is the agent-side counterpart, built before generation and
verified after it. The two compose naturally: a provenance-tagged
specification field is an ideal anchor for selecting the corresponding
subtree in an editable assembly view, and we view intent-level provenance
plus structure-level editing as the full-stack answer to the
``regenerate-everything'' pathology.

\paragraph{Honest boundaries.}
The verifier's reach ends where stated intent ends: unstated errors,
wrong topology never discussed with the user, and semantics beyond the
feature vocabulary remain the province of offline similarity metrics and,
eventually, human inspection. Under lawful non-uniqueness (underspecified
prompts with many valid resolutions), dataset-grounded Chamfer penalizes
valid confirmations; we therefore separate dataset-referenced scores from
specification-referenced ones rather than average over both.

\paragraph{On the number of agents.}
\label{sec:discussion:agents}
Our system has three agents---Clarifier, Spec Manager, and Coder---plus a
deterministic verifier that is a \emph{tool}, not an agent: it has no LLM, no
objectives, and no action choice. We deliberately stop at three. The
Clarifier and Coder are inherited from ProCAD and carry well-understood,
separately trained responsibilities; the Spec Manager is the single new
agent because it is the one component that genuinely needs state and
user-facing decisions---what to ask, which value to trust, where to route a
failure. The verifier stays a tool for a principled reason: a
specification--program checker is a total function, and naming it an agent
would invite exactly the ``multi-agent as prompt engineering'' criticism
that undermines the method's claim to verifiability. CADDesigner achieves a
comparable feature set inside one general-purpose agent \cite{caddesigner};
our ablation- and unit-test-friendly decomposition is an alternative whose
trade-off---more components, each with a smaller, mechanically checkable
contract---we argue is the right one when correctness, not fluency, is the
requirement. Three agents are also not a floor: a natural collapse is that
the Clarifier is eventually absorbed into the Spec Manager, whose
unresolved-field machinery already subsumes question asking, leaving a
two-agent system with the same guarantees. We leave that consolidation to
future work and report the three-agent form here because it isolates the
Spec Manager's contribution against the ProCAD baseline cleanly.

\paragraph{Conclusion.}
We replaced the last free-text hand-off in clarify-then-generate text-to-CAD
pipelines with a provenance-tracked specification governed by a deterministic
kernel, and grounded verification in what the user actually confirmed rather
than in ground-truth geometry or unanchored self-inspection. Information
loss, invention, and carried-over errors---invisible in free-text
pipelines---become measurable, preventable, and repairable at deployment
time, at question budgets matched to baselines, and reproducible on open
7B models throughout.

\paragraph{Reproducibility.}
All code (pipeline, deterministic kernel with offline unit tests, verifier,
distillation, and analysis scripts) and the SFT recipe are released; the
evaluation reuses the public ProCAD dataset and checkpoints.
